\begin{abstract}
针对光伏不确定性、负载波动、储能跨时段耦合以及外网电价变化共同作用下的微网购电问题，本文建立一套由\textbf{确定性线性规划---风险修正日前计划---多时刻滚动调整---因果价格预测}逐层扩展的统一优化框架。首先以 10 min 为基本时间尺度，将购电量、充放电量和储能电量统一到能量平衡关系中，以线性规划最小化购电费用；其次对全年负载与光伏数据实施逐日递推（walk-forward）预测，并依据紧急购电价为正常电价 5 倍这一不对称损失结构，引入过去 28 d 净负荷残差的 0.8 分位数作为风险安全裕度；随后利用 0:00、6:00、12:00、18:00 的光伏预报构建滚动时域控制，在外部预报与历史预报之间按滚动 MAE 选择，并根据当天已实现负载修正剩余时段预测；最后针对问题四，避免直接使用当日未来真实电价造成信息泄露，采用 7 d 滞后同时间隔价格作为可实施的因果预测，并把``完整已知当日电价''仅作为完全信息下界（oracle）。

计算结果表明：问题一优化后全天购电量为 \textbf{59,482.699 kWh}，购电费为 \textbf{35,126.949 元}，较无储能净负荷直购对照方案降低 \textbf{26.90\%}。问题二在 2025-02-01 至 2025-12-31 的总费用为 \textbf{14,916,093.656 元}，紧急购电量为 \textbf{367,486.925 kWh}。引入日内滚动预报后，问题三总费用降至 \textbf{14,679,141.979 元}，较问题二下降 \textbf{1.59\%}，紧急购电量下降 \textbf{28.98\%}；其中 6:00、12:00 的外部光伏预报在滚动择优中选择率均为 100\%，说明这两个时刻的更新最值得纳入调整。波动电价下，严格因果的 Q4-2 与 Q4-3 总费用分别为 \textbf{15,774,011.240 元}和 \textbf{15,507,241.002 元}，滚动调整仍节省 \textbf{1.691\%}；7 d 滞后价格预测的 MAE 为 \textbf{0.047151 元/kWh}，两种因果策略距完全价格信息 oracle 仅约 \textbf{1.09\%/1.10\%}。对风险分位数 0.7/0.8/0.9 及问题三下调结算口径的敏感性分析表明，主结论在当前数据与题设场景下稳定。

\textbf{关键词：} 微网；储能调度；线性规划；滚动时域优化；风险分位数；光伏预测；因果价格预测
\end{abstract}
